% =====================================================================
\chapter{Comprensión de los datos}
\label{cap:datos}

Este capítulo corresponde a la segunda fase de CRISP-DM: qué datos hay, de dónde vienen y qué
problemas tienen. Hay dos fuentes, las noticias y los precios, y cada una se trata en su
carpeta del proyecto (\carpeta{2_comprension_de_los_datos/noticias} y
\carpeta{2_comprension_de_los_datos/mercado}).

% ---------------------------------------------------------------------
\section{Las noticias}
\label{sec:recoleccion}

No existe un archivo público y ordenado de las notas de economía de los portales argentinos, así
que el corpus se construyó con \textit{web scraping}: la extracción automática de información de
páginas web. Se usaron dos caminos, según la disponibilidad de cada portal:

\begin{itemize}
  \item \textbf{Internet Archive.} Su \textit{Wayback Machine} guarda capturas de millones de
        páginas, y su API permite listar las capturas de un sitio en un rango de fechas. Así se
        reconstruye una ventana más larga que la que exponen los propios portales, a costa de
        depender de qué capturó un tercero \citep{garg2025wayback}. Se usó para Ámbito Financiero,
        Clarín (Economía), Infobae, La Nación y TN.
  \item \textbf{Mapas del sitio.} Los portales publican archivos \texttt{sitemap.xml} para los
        buscadores, que listan sus notas recientes. Se usaron para Clarín, El Economista e
        iProfesional.
\end{itemize}

De cada página se separó el texto de la nota del «relleno» (menús, publicidad, firmas) con una
estrategia en capas: los datos estructurados JSON-LD, las metaetiquetas Open Graph, la librería
\texttt{trafilatura} \citep{barbaresi2021} y, como último recurso, reglas sobre el HTML. Cada nota
se guardó con su portal, fecha y hora de publicación, título, bajada y contenido. El corpus final
tiene \textbf{\dato{corpus.notas} notas} de \dato{corpus.portales} portales, recolectadas en
\dato{corpus.fuentes} fuentes (tabla~\ref{tab:corpus}).

\begin{table}[H]
\centering\small
\begin{tabular}{@{}llr@{}}
\toprule
\textbf{Portal} & \textbf{Método de recolección} & \textbf{Notas} \\
\midrule
\tabla{corpus}
\bottomrule
\end{tabular}
\caption{Corpus de noticias por fuente, después de la depuración. Fuente: elaboración propia
(\carpeta{2_comprension_de_los_datos/noticias}).}
\label{tab:corpus}
\end{table}

\textbf{Auditoría y depuración.} Las herramientas de extracción de contenido fallan con frecuencia
en el relleno y en los metadatos \citep{foley2026wcxb}, y los duplicados distorsionan cualquier
modelo que se entrene con ese texto \citep{lee2021dedup}, por lo que un corpus web tiene que
auditarse antes de usarse \citep{dodge2021c4}. La auditoría encontró cinco problemas, cada uno con
su solución (tabla~\ref{tab:depuracion}).

\begin{table}[H]
\centering\small
\begin{tabularx}{\linewidth}{@{}L{3.2cm}YY@{}}
\toprule
\textbf{Problema} & \textbf{Qué se encontró} & \textbf{Solución} \\
\midrule
Duplicados exactos & 205 notas de Clarín descargadas dos veces, por la Wayback Machine y por el mapa del sitio & Una copia por URL \\
Publicación cruzada & 208 grupos de título repetido; 33 grupos (68 notas) a menos de 48 horas & Verificación con el modelo de lenguaje Qwen3: 14 republicaciones reales, eliminadas \\
Metadatos faltantes & 194 notas de El Economista sin fecha, por un error de lectura del JSON-LD & Lectura tolerante; 192 notas recuperaron su fecha \\
Contenido que no es noticia & Páginas índice, de suscripción y notas fechadas por su última modificación & Se eliminaron \\
Plantillas & Notas como «Dólar hoy», con texto fijo y un sentimiento constante que sesgaba el promedio diario & Se eliminaron 419 notas (2,9\%): el corpus pasó de 14.431 a \dato{corpus.notas} \\
\bottomrule
\end{tabularx}
\caption{Problemas del corpus y cómo se resolvieron. Fuente: elaboración propia a partir de la
auditoría del corpus (\carpeta{2_comprension_de_los_datos/noticias/documentacion}).}
\label{tab:depuracion}
\end{table}

\textbf{Límites de la fuente.} La cobertura depende de qué capturó la Wayback Machine y de cuánto
tiempo guardan los portales sus mapas del sitio. Por eso la \textit{cantidad} de notas de un día
no es una medida confiable de cuánto se habló de economía ese día: depende también de la
recolección. Abril de 2026, por ejemplo, no tiene notas de iProfesional, el portal más grande. Esta
es la razón por la que la cantidad de notas no entra en los modelos: se usa solo el sentimiento,
que es un promedio y depende mucho menos de cuántas notas hay. Las notas pertenecen a sus medios y
se usaron con fines académicos: el análisis trabaja con medidas agregadas y no reproduce las
notas, salvo oraciones sueltas como ejemplo.

% ---------------------------------------------------------------------
\section{El mercado}
\label{sec:panel}

Los precios diarios (apertura, máximo, mínimo, cierre y volumen) se descargaron de Yahoo Finance
con \texttt{yfinance}. Se partió de \dato{mercado.candidatas} acciones candidatas del Panel General
y del Panel Líderes de BYMA, y se exigieron dos condiciones en los tres años de datos
(\dato{mercado.ruedas_3anos} ruedas): cobertura completa del calendario y al menos 90\% de días
con volumen real, porque un día sin operaciones no es un faltante sino una fila con el precio
pegado. Pasaron \textbf{\dato{mercado.acciones} acciones} de \dato{mercado.sectores} sectores,
incluidas las \dato{mercado.en_merval} del S\&P~Merval. En la depuración se quitaron las filas que
Yahoo agrega en los feriados (se usó como calendario de BYMA la lista de fechas del índice) y se
corrigieron dos eventos corporativos mal ajustados: el \textit{split} de SEMI de diciembre de 2024
y el dividendo en acciones de AGRO de noviembre de 2023, que mostraba una caída falsa de 88\%.

\textbf{La ventana de estudio.} Las noticias empiezan en la práctica en mayo de 2025, porque tanto
los mapas del sitio como las capturas de la Wayback Machine se concentran en los meses más
recientes. La ventana común entre noticias y precios va del 05/05/2025 al 30/04/2026:
\textbf{\dato{dataset.ruedas} ruedas}. De las \dato{mercado.acciones} acciones, el modelo usa las
\dato{dataset.empresas} que tienen noticias en suficientes ruedas como para medir su sentimiento
propio (sección~\ref{sec:propio}).
